% Кодировка и язык
\usepackage[T2A]{fontenc} % поддержка кириллицы
\usepackage[utf8]{inputenc} % кодировка исходного текста
\usepackage[english,russian]{babel} % переключение языков

% Геометрия страницы и графика
\usepackage[left=2.5cm, right=2cm, top=2cm, bottom=2cm]{geometry} % поля страницы
\usepackage{graphicx} % подключение графики
\usepackage{pdfpages} % вставка pdf-страниц

% Таблицы
\usepackage{array} % расширенные возможности для работы с таблицами
\usepackage{tabularx} % автоматический подбор ширины столбцов
\usepackage{dcolumn} % выравнивание чисел по разделителю

% Математика
\usepackage{bm} % полужирное начертание для математических символов
\usepackage{amsmath} % дополнительные математические возможности
\usepackage{amssymb} % дополнительные математические символы

% Библиография и ссылки
\usepackage{cite} % поддержка цитирования
\usepackage[hyperfootnotes=false,hidelinks]{hyperref} % создание гиперссылок
\providecommand{\selectlanguageifdefined}[1]{%
  \expandafter\ifx\csname date#1\endcsname\relax
  \else\selectlanguage{#1}\fi
}
\providecommand*{\BibUrl}[1]{\url{#1}}
\providecommand{\BibAnnote}[1]{}
\providecommand*{\BibEmph}[1]{#1}
\ProvideTextCommandDefault{\cyrdash}{\iflanguage{russian}{\hbox
  to.8em{--\hss--}}{\textemdash}}
\providecommand*{\BibDash}{}
\providecommand{\newblock}{\ignorespaces}
\makeatletter
\newcommand{\fullciteentry}[1]{%
  \@ifundefined{fullcite@#1}{%
    \textbf{[Источник #1 не найден; обновите полные библиографические записи]}%
  }{%
    \csname fullcite@#1\endcsname
  }%
}
\newcommand{\printfootcitelist}[1]{%
  \begingroup
  \def\href##1##2{##2}%
  \def\url##1{\nolinkurl{##1}}%
  \def\BibUrl##1{\nolinkurl{##1}}%
  \def\footcite@separator{}%
  \@for\footcite@key:=#1\do{%
    \footcite@separator\expandafter\fullciteentry\expandafter{\footcite@key}%
    \def\footcite@separator{\par\medskip}%
  }%
  \endgroup
}
\DeclareRobustCommand{\footcite}[1]{\footnote{\nocite{#1}\printfootcitelist{#1}}\selectlanguage{russian}}
\makeatother
\InputIfFileExists{biblio/footcite-entries.tex}{}{}

% Прочее
\usepackage{color} % работа с цветом
\usepackage{epstopdf} % конвертация eps в pdf
\usepackage{multirow} % объединение ячеек таблиц по вертикали
\usepackage{afterpage} % вставка материала после текущей страницы
\usepackage[font={normal}]{caption} % настройка подписей к рисункам и таблицам
\usepackage[onehalfspacing]{setspace} % полуторный интервал
\usepackage{fancyhdr} % установка колонтитулов
\usepackage{listings} % поддержка вставки исходного кода

% Установка шрифта Times New Roman
\renewcommand{\rmdefault}{ftm}

% Настройка текста
\usepackage{setspace}
\setstretch{1.5} % промежуточный интервал
\setlength{\parindent}{1.5cm}


% Создание нового типа столбца для выравнивания содержимого по центру
\newcommand{\PreserveBackslash}[1]{\let\temp=\\#1\let\\=\temp}
\newcolumntype{C}[1]{>{\PreserveBackslash\centering}p{#1}}
\newcolumntype{L}[1]{>{\raggedright\arraybackslash}p{#1}}
\newcolumntype{Y}{>{\raggedright\arraybackslash}X}

% Настройка стиля страницы
\pagestyle{fancy}      % Использование стиля "fancy" для оформления страниц
\fancyhf{}              % Очистка текущих значений колонтитулов
\fancyfoot[C]{\thepage} % Установка номера страницы в нижнем колонтитуле по центру
\renewcommand{\headrulewidth}{0pt} % Удаление разделительной линии в верхнем колонтитуле

% Настройка подписей к изображениям и таблицам
\captionsetup{format=hang,labelsep=period}
\DeclareCaptionLabelSeparator{tablelinebreak}{\par}
\captionsetup[table]{
  position=top,
  format=plain,
  labelsep=tablelinebreak,
  justification=centering,
  singlelinecheck=false
}

% Использование полужирного начертания для векторов
\let\vec=\mathbf

% Установка глубины оглавления
\setcounter{tocdepth}{2}

% Указание папки для поиска изображений
\graphicspath{{images/}}

% Установка стилей страницы и главы
\pagestyle{footcenter}
\chapterpagestyle{footcenter}
